% ------------------------------------------------------------
% 1. DOCUMENTO E IDIOMA
% ------------------------------------------------------------

\usepackage[utf8]{inputenc}
\usepackage[T1]{fontenc}
\usepackage[english,spanish,es-tabla]{babel}

\decimalpoint

% ------------------------------------------------------------
% 2. TIPOGRAFÍA
% ------------------------------------------------------------

\usepackage{newtxtext}
\usepackage{microtype}

% ------------------------------------------------------------
% 3. CONFIGURACIÓN DE PÁGINA - APA 7
% ------------------------------------------------------------

\usepackage[
    letterpaper,
    margin=1in,
    headheight=15pt
]{geometry}

\usepackage{setspace}
\doublespacing

\setlength{\parindent}{0.5in}
\setlength{\parskip}{0pt}
\setlength{\emergencystretch}{3em}

% ------------------------------------------------------------
% 4. MATEMÁTICAS Y FÍSICA
% ------------------------------------------------------------

\usepackage{amsmath}
\usepackage{amssymb}
\usepackage{mathtools}
\usepackage{bm}

% Unidades físicas
\usepackage{siunitx}

\sisetup{
    output-decimal-marker = {.},
    group-separator = {\,},
    group-minimum-digits = 5,
    per-mode = symbol,
    separate-uncertainty = true
}

% ------------------------------------------------------------
% 5. FIGURAS
% ------------------------------------------------------------

\usepackage{graphicx}
\usepackage{float}
\usepackage{caption}
\usepackage{xcolor}

% Directorios de imágenes
\graphicspath{
    {imagenes/}
    {imagenes/figuras/}
}

% Configuración de títulos de figuras y tablas

\captionsetup{
    font=normalsize,
    labelfont=bf,
    textfont=it,
    labelsep=newline,
    justification=raggedright,
    singlelinecheck=false,
    skip=0.25\baselineskip
}

% ------------------------------------------------------------
% 6. TABLAS
% ------------------------------------------------------------

\usepackage{booktabs}
\usepackage{array}

% ------------------------------------------------------------
% 7. CÓDIGO PYTHON
% ------------------------------------------------------------

\usepackage{listings}
\usepackage{upquote}

\lstdefinestyle{pythonAPA}{
    language=Python,
    basicstyle=\ttfamily\small,
    numbers=left,
    numberstyle=\tiny,
    numbersep=8pt,
    stepnumber=1,
    showstringspaces=false,
    breaklines=true,
    breakatwhitespace=false,
    keepspaces=true,
    columns=fullflexible,
    tabsize=4,
    frame=single,
    rulecolor=\color{black},
    captionpos=b,
    inputencoding=utf8,
    literate=
        {á}{{\'a}}1
        {é}{{\'e}}1
        {í}{{\'i}}1
        {ó}{{\'o}}1
        {ú}{{\'u}}1
        {Á}{{\'A}}1
        {É}{{\'E}}1
        {Í}{{\'I}}1
        {Ó}{{\'O}}1
        {Ú}{{\'U}}1
        {ñ}{{\~n}}1
        {Ñ}{{\~N}}1
        {ü}{{\"u}}1
        {Ü}{{\"U}}1
        {×}{{$\times$}}1
        {°}{{$^\circ$}}1
        {→}{{$\rightarrow$}}1
        {←}{{$\leftarrow$}}1
}

% ------------------------------------------------------------
% 8. LISTAS
% ------------------------------------------------------------

\usepackage{enumitem}

\setlist{
    topsep=0pt,
    itemsep=0pt,
    parsep=0pt,
    partopsep=0pt,
    leftmargin=0.5in
}

% ------------------------------------------------------------
% 9. ENCABEZADO Y NUMERACIÓN DE PÁGINAS
% ------------------------------------------------------------

\usepackage{fancyhdr}

\pagestyle{fancy}
\fancyhf{}

\fancyhead[R]{\thepage}

\renewcommand{\headrulewidth}{0pt}
\renewcommand{\footrulewidth}{0pt}

\fancypagestyle{plain}{
    \fancyhf{}
    \fancyhead[R]{\thepage}
    \renewcommand{\headrulewidth}{0pt}
}

% ------------------------------------------------------------
% 10. FORMATO DE TÍTULOS - APA 7
% ------------------------------------------------------------

\usepackage{titlesec}

% Nivel 1
\titleformat{\chapter}[block]
    {\normalfont\normalsize\bfseries\centering}
    {}{0pt}{}

\titlespacing*{\chapter}
    {0pt}{0pt}{\baselineskip}

% Nivel 2
\titleformat{\section}[block]
    {\normalfont\normalsize\bfseries}
    {}{0pt}{}

\titlespacing*{\section}
    {0pt}{\baselineskip}{0pt}

% Nivel 3
\titleformat{\subsection}[block]
    {\normalfont\normalsize\bfseries\itshape}
    {}{0pt}{}

\titlespacing*{\subsection}
    {0pt}{\baselineskip}{0pt}

% Nivel 4
\titleformat{\subsubsection}[runin]
    {\normalfont\normalsize\bfseries}
    {}{0pt}{}[.]

\titlespacing*{\subsubsection}
    {0pt}{\baselineskip}{0.5em}

% Nivel 5
\titleformat{\paragraph}[runin]
    {\normalfont\normalsize\bfseries\itshape}
    {}{0pt}{}[.]

\titlespacing*{\paragraph}
    {0pt}{\baselineskip}{0.5em}

% APA: encabezados sin numeración
\setcounter{secnumdepth}{-1}
\setcounter{tocdepth}{2}

% ------------------------------------------------------------
% 11. BIBLIOGRAFÍA - APA 7
% ------------------------------------------------------------

\usepackage{csquotes}

\usepackage[
    backend=biber,
    style=apa,
    sorting=nyt,
    sortcites=true,
    uniquelist=false
]{biblatex}

\DeclareLanguageMapping{spanish}{spanish-apa}

% Sangría francesa
\setlength{\bibhang}{0.5in}
\setlength{\bibitemsep}{0pt}

% ------------------------------------------------------------
% 12. HIPERVÍNCULOS Y REFERENCIAS CRUZADAS
% ------------------------------------------------------------

\usepackage[hidelinks]{hyperref}

\usepackage[
    nameinlink,
    noabbrev,
    capitalise
]{cleveref}

\crefname{figure}{figura}{figuras}
\Crefname{figure}{Figura}{Figuras}

\crefname{table}{tabla}{tablas}
\Crefname{table}{Tabla}{Tablas}

\crefname{equation}{ecuación}{ecuaciones}
\Crefname{equation}{Ecuación}{Ecuaciones}

% ------------------------------------------------------------
% 13. DATOS INSTITUCIONALES
% ------------------------------------------------------------

\newcommand{\universidad}{Universidad de Nariño}

\newcommand{\facultad}{
    Facultad de Ciencias Exactas y Naturales
}

\newcommand{\programa}{Programa de Física}

\newcommand{\ciudad}{San Juan de Pasto, Nariño}

% ------------------------------------------------------------
% 14. DATOS DEL INFORME
% ------------------------------------------------------------

\newcommand{\tituloInforme}{
    Modelación computacional de un acelerador lineal
}

\newcommand{\autorInforme}{
    Erick García\\
    Gabriela Herrera\\
    Juan Esteban Muñoz
}

\newcommand{\docenteInforme}{
    Mg. Jaime Betancourth
}

\newcommand{\asignaturaInforme}{
    Herramientas Computacionales para la Física II
}

\newcommand{\fechaInforme}{\today}


